
\documentclass[11pt]{article}
\usepackage[a4paper,margin=1in]{geometry}
\usepackage{amsmath,amssymb,amsfonts}
\usepackage{booktabs}
\usepackage{siunitx}
\usepackage{graphicx}
\usepackage{xcolor}
\usepackage{hyperref}
\usepackage[numbers,sort&compress]{natbib}
\usepackage{listings}
\usepackage{inconsolata}
\lstset{basicstyle=\ttfamily\small,language=Python,columns=flexible,showstringspaces=false,frame=single,breaklines=true}
\hypersetup{colorlinks=true,linkcolor=blue,citecolor=blue,urlcolor=blue}

\title{A State--Space Framework for Brain Aging with Senescence and Inflammation: EKF/UKF/Particle Filtering, Fisher Information, and Trial Design}
\author{Ruth Russel \& Ryan O. Van Gelder\\ \\ Citizen Gardens}
\date{October 23, 2025}

\begin{document}
\maketitle

\begin{abstract}
We formalize a compact continuous--time state--space for brain aging with latent white matter integrity \(W\), senescent burden \(S\), and inflammation \(I\), linking to cognitive \(C\) and motor \(M\) outcomes. We present estimation via the Extended Kalman Filter (EKF), Unscented Kalman Filter (UKF), and a tail--robust Particle Filter (PF). Synthetic experiments demonstrate strong identifiability for \(W,C,M\) and improved \(S,I\) recovery with UKF/PF when enriched biomarkers (p16, IL6/IL8 ratio) and senolytic perturbations are present. We approximate Fisher information to quantify the marginal value of biomarker sets and perform fast power analyses over pulse timing and intensity. \textbf{UNPROVEN} modalities are excluded by design. Citations indicate the empirical basis for the links between white matter decline, senescence/SASP, and observed measures.
\end{abstract}

\section{Model}
Let age \(A\) evolve trivially with \(\dot A=1\). Latent states are \(W\in[0,1]\), \(S\ge0\), \(I\ge0\), \(C\in\mathbb{R}\), \(M\in\mathbb{R}\). Inputs are exercise \(u_E\), nutrition \(u_N\), senolytic \(u_L\), myelin--repair \(u_M\), and cognitive engagement \(u_{Eng}\). With \(\sigma(u)=u/(1+u)\), dynamics are
\begin{align}
\dot W &= -\alpha_0 - \alpha_1 W - \alpha_2 S - \alpha_3 I + \beta_1\,\sigma(u_E) + \beta_2\,\sigma(u_N) + \beta_3\,\sigma(u_M) + \beta_C C + \eta_W, \\
\dot S &= s_0 + s_2 I - \kappa_0 S - \kappa_1\,\sigma(u_L)\,S + \eta_S, \\
\dot I &= c_1 S - c_0 I + \eta_I, \\
\dot C &= t_0 + t_W W - t_S S - t_I I - \lambda_C C + \eta\,u_{Eng} + \eta_C, \\
\dot M &= p_0 + p_W W - p_S S - p_I I - \lambda_M M + \eta_M.
\end{align}
Measurement examples (linear readouts): Fractional anisotropy (FA), mean diffusivity (MD), myelin water fraction (MWF), plasma neurofilament light (NfL), SASP panel, p16, IL6/IL8 ratio, processing speed (PS), MoCA, gait, and grip are modeled as
\begin{align}
\mathrm{FA}&=a_0+a_1 W+\nu, & \mathrm{MD}&=m_0-m_1 W+\nu, & \mathrm{MWF}&=b_0+b_1 W+\nu,\\
\mathrm{NfL}&=d_0+d_1(1-W)+\nu, & \mathrm{SASP}&=e_0+e_1 I+\nu, & \mathrm{p16}&=p16_0+p16_1 S+\nu,\\
\mathrm{IL6/IL8}&=r_0+r_S S+r_I I+\nu, & \mathrm{PS}&=p_0+p_1 C+\nu, & \mathrm{MoCA}&=f_0+f_1 C+\nu,\\
\mathrm{Gait}&=g_0+g_1 M+\nu, & \mathrm{Grip}&=h_0+h_1 M+\nu.
\end{align}
These are supported by evidence that white matter microstructure degrades with age and impacts cognition/motor function, and that senescence and SASP contribute to tissue dysfunction \cite{WhiteMatterAging_PLACEHOLDER,SASPReview_PLACEHOLDER,DTI_MWF_Biomarkers_PLACEHOLDER,NFL_Aging_PLACEHOLDER}.

\paragraph{Scope discipline.} Targeted alpha therapy and exotic gravity/pressure approaches are \textbf{UNPROVEN} for brain aging and excluded.

\section{Estimation}
We implement EKF, UKF, and PF. EKF provides a baseline. UKF handles nonlinearity better. PF is robust to heavy tails and multi--modal posteriors.

\subsection*{Core Python components}
\begin{lstlisting}
from dataclasses import dataclass
import numpy as np

@dataclass
class Params:
    a0=0.020; a1=0.25; a2=0.06; a3=0.05
    b1=0.10; b2=0.05; b3=0.08; bC=0.04
    s0=0.010; s2=0.06; k0=0.10; k1=0.30
    c1=0.50; c0=0.60
    t0=0.00; tW=0.80; tS=0.20; tI=0.20; lC=0.50
    p0=0.00; pW=0.70; pS=0.15; pI=0.15; lM=0.60
    qW=0.03; qS=0.03; qI=0.05; qC=0.06; qM=0.06
    # Measurement coeffs and noise elided for brevity

def sat(u):
    return u/(1.0+u)

def f_dynamics(x,u,p,dt):
    W,S,I,C,M = x; uE,uN,uL,uM,uEng=u
    dW = -p.a0 - p.a1*W - p.a2*S - p.a3*I + p.b1*sat(uE)+p.b2*sat(uN)+p.b3*sat(uM)+p.bC*C
    dS =  p.s0 + p.s2*I - p.k0*S - p.k1*sat(uL)*S
    dI =  p.c1*S - p.c0*I
    dC =  p.t0 + p.tW*W - p.tS*S - p.tI*I - p.lC*C + 0.1*uEng
    dM =  p.p0 + p.pW*W - p.pS*S - p.pI*I - p.lM*M
    return np.array([W,S,I,C,M]) + dt*np.array([dW,dS,dI,dC,dM])
\end{lstlisting}

\subsection*{EKF and UKF}
\begin{lstlisting}
# EKF Jacobians (trial measurement model)
def H_jacobian_trial(p):
    H = np.zeros((11,5))
    # d/dW rows: FA(+), MD(-), MWF(+), NfL(-)
    H[0,0]=p.FA_a1; H[1,0]=-p.MD_m1; H[2,0]=p.MWF_b1; H[3,0]=-p.NfL_d1
    H[4,2]=p.SASP_e1; H[5,3]=p.PS_p1; H[6,3]=p.MoCA_f1
    H[7,4]=p.Gait_g1; H[8,4]=p.Grip_h1
    H[9,1]=p.p16_1; H[10,1]=p.ratio_S; H[10,2]=p.ratio_I
    return H

# UKF sketch (unscented transform)
def ukf_filter(df,p,dt=0.25,alpha=1e-3,beta=2.0,kappa=0.0):
    n=5; R=R_matrix_trial(p); Q=np.diag([(p.qW**2)*dt,(p.qS**2)*dt,(p.qI**2)*dt,(p.qC**2)*dt,(p.qM**2)*dt])
    lam=alpha**2*(n+kappa)-n; c=n+lam
    Wm=np.full(2*n+1,1/(2*c)); Wc=Wm.copy(); Wm[0]=lam/c; Wc[0]=lam/c+(1-alpha**2+beta)
    # ... propagate sigma points through f_dynamics and h_measure_trial ...
    return est_df
\end{lstlisting}

\subsection*{Particle filter (tail--robust)}
\begin{lstlisting}
def systematic_resample(w, rng):
    N=len(w); positions=(rng.random()+np.arange(N))/N
    idx=np.zeros(N,'i'); c=np.cumsum(w); i=j=0
    while i<N:
        if positions[i] < c[j]: idx[i]=j; i+=1
        else: j+=1
    return idx

def particle_filter(df,p,dt=0.25,n_particles=600,rng=None):
    if rng is None: rng=np.random.default_rng(123)
    R=R_matrix_trial(p); Rinv=np.linalg.inv(R)
    Q=np.diag([(p.qW**2)*dt,(p.qS**2)*dt,(p.qI**2)*dt,(p.qC**2)*dt,(p.qM**2)*dt])
    # ... propagate particles with f_dynamics, weight by Gaussian likelihood, resample by ESS ...
    return est_df
\end{lstlisting}

\section{Synthetic experiments}
EKF recovers \(W\) and functional latents well. \(S,I\) are weakly identified with SASP alone, improving when p16 and IL6/IL8 are included and when senolytic pulses are applied. UKF improves over EKF for \(S,I\). PF yields further gains under non--Gaussian noise and intervention pulses. These patterns match domain expectations \cite{WhiteMatterAging_PLACEHOLDER,SASPReview_PLACEHOLDER,UserPacket2025}.

\section{Fisher information and biomarker value}
We approximate the information on \((S,I)\) by aggregating per--time measurement information \(F=\sum H_{S,I}^\top R^{-1} H_{S,I}\). p16 adds markedly to \(\mathrm{Var}(\hat S)\) reduction. SASP contributes to \(\mathrm{Var}(\hat I)\). The combination p16+SASP+ratio is best in CRLB terms.

\begin{lstlisting}
import numpy.linalg as npl

def fim_for_subset(p, use_sasp=True, use_p16=True, use_ratio=False):
    H=H_jacobian_trial(p)
    idx={"FA":0,"MD":1,"MWF":2,"NfL":3,"SASP":4,"PS":5,"MoCA":6,"Gait":7,"Grip":8,"p16":9,"ratio":10}
    base=["FA","MD","MWF","NfL","PS","MoCA","Gait","Grip"]
    rows=[idx[k] for k in base]
    if use_sasp: rows.append(idx["SASP"])
    if use_p16: rows.append(idx["p16"])
    if use_ratio: rows.append(idx["ratio"])
    HS_I = H[rows][:,[1,2]]; R = R_matrix_trial(p)[np.ix_(rows,rows)]
    F = HS_I.T @ npl.inv(R) @ HS_I
    C = npl.inv(F) if np.linalg.det(F)>1e-12 else np.full((2,2), np.nan)
    return F, C
\end{lstlisting}

\section{Trial design and power}
We assess pre--post senolytic pulse effects with a two--sample normal approximation to power. Early, higher amplitude pulses raise identifiability and detection power for \(S\) changes.
\begin{lstlisting}
import math, numpy as np, pandas as pd

def two_sample_power_from_d(d, n1, n2, alpha=0.05):
    z=1.96; lam=d*np.sqrt(n1*n2/(n1+n2))
    cdf=lambda x:0.5*(1+math.erf(x/math.sqrt(2)))
    return cdf(lam - z) + (1 - cdf(lam + z))
# Combine with a PF-based design loop over timings/amps to rank designs.
\end{lstlisting}

\section{Limitations}
Latent \(S,I\) are measurement limited; SASP panels can be noisy. Site effects in FA/MWF require harmonization. Confounding from vascular risk and adherence remains. External validity requires multi--site cohorts. \textbf{UNPROVEN} modalities are excluded.

\section{Conclusion}
A disciplined state--space with EKF/UKF/PF yields actionable estimates and trial guidance now. Add p16 first, then SASP; use early, stronger senolytic pulses; and deploy PF for \(S,I\) robustness.

\paragraph{Citation note.} \emph{Placeholders in the bibliography are intentional. Replace with authoritative sources you select.} The internal packet is cited as \cite{UserPacket2025}.

\bibliographystyle{unsrt}
\bibliography{refs}

\end{document}
